\documentclass[11pt, a4paper]{article}
\usepackage[a4paper, margin=2.5cm]{geometry}
\usepackage{fontspec}
\usepackage[vietnamese, english, bidi=basic, provide=*]{babel}
\babelprovide[import, onchar=ids fonts]{english}
\babelfont{rm}{Noto Sans}
\usepackage{enumitem}
\setlist[itemize]{label=-}
\usepackage{amsmath, amssymb}
\usepackage{booktabs}
\usepackage{hyperref}

\title{\textbf{BÁO CÁO ĐỀ TÀI 10: ĐẶT CAMERA GIÁM SÁT GIAO THÔNG (BÀI TOÁN PHỦ TẬP HỢP -- SET COVER)}\\
\large Phần A: Cơ sở lý thuyết, Thuật toán Chính xác và Simulated Annealing}
\author{Nhóm 09}
\date{}

\begin{document}

\maketitle

\section{Phát biểu bài toán}

\subsection{Bối cảnh và mô tả hình thức}
Cơ quan quản lý giao thông đô thị cần lắp đặt hệ thống camera để giám sát toàn bộ các nút giao trọng điểm của một quận. Các vị trí lắp đặt ứng viên (cột đèn, giá long môn, nóc tòa nhà) có chi phí đầu tư khác nhau và mỗi vị trí quan sát được một tập hợp các nút giao lân cận. Mục tiêu là lựa chọn một tập con các vị trí đặt camera có tổng chi phí nhỏ nhất sao cho mọi nút giao đều được ít nhất một camera quan sát.

Trong quy ước của thư viện chuẩn OR-Library, các nút giao tương ứng với các hàng và các vị trí đặt camera tương ứng với các cột của ma trận liên thuộc. Đầu vào bài toán gồm:
\begin{itemize}
    \item Tập vũ trụ $U = \{1, 2, \dots, m\}$ gồm $m$ nút giao (hàng) cần được giám sát.
    \item Họ gồm $n$ tập con $\mathcal{S} = \{S_1, S_2, \dots, S_n\}$ tương ứng với $n$ vị trí ứng viên (cột), trong đó $S_j \subseteq U$ là tập các nút giao được quan sát bởi vị trí $j$, thỏa mãn điều kiện khả thi $\bigcup_{j=1}^{n} S_j = U$.
    \item Mỗi vị trí $j \in \{1, \dots, n\}$ có chi phí lắp đặt dương $c_j > 0$.
\end{itemize}

\subsection{Mô hình Quy hoạch nguyên tuyến tính (ILP) và không gian tìm kiếm}
Đặt biến quyết định nhị phân $x_j \in \{0, 1\}$ cho mỗi vị trí $j \in \{1, \dots, n\}$, trong đó $x_j = 1$ nếu chọn đặt camera tại vị trí $j$ và $x_j = 0$ nếu ngược lại. Ký hiệu phần tử của ma trận phủ $A = (a_{ij})_{m \times n}$ với $a_{ij} = 1$ nếu $i \in S_j$ và $a_{ij} = 0$ nếu $i \notin S_j$. Mô hình ILP của bài toán được viết như sau:
\begin{align}
\text{Tối thiểu hóa:} \quad & Z = \sum_{j=1}^{n} c_j x_j \\
\text{Ràng buộc:} \quad & \sum_{j : i \in S_j} x_j \ge 1, \quad \forall i \in \{1, 2, \dots, m\} \\
& x_j \in \{0, 1\}, \quad \forall j \in \{1, 2, \dots, n\}
\end{align}

\textbf{Kích thước không gian tìm kiếm:} Do mỗi vị trí $j$ có hai trạng thái chọn hoặc không chọn, không gian tìm kiếm toàn cục gồm $2^n$ cấu hình nhị phân khác nhau.

\subsection{Ví dụ minh họa nhỏ}
Xét thực thể ví dụ của Đề 10 gồm $m = 7$ nút giao và $n = 6$ vị trí đặt camera:
\begin{itemize}
    \item Chi phí các vị trí: $c = (7, 5, 9, 4, 3, 7)$.
    \item Tập nút giao phủ được theo từng cột: $S_1 = \{1, 4, 7\}$, $S_2 = \{5, 6\}$, $S_3 = \{2, 3, 5, 7\}$, $S_4 = \{2, 3, 6, 7\}$, $S_5 = \{4, 6\}$, $S_6 = \{1, 3, 6, 7\}$.
\end{itemize}
Duyệt toàn bộ $2^6 = 64$ cấu hình cho thấy nghiệm tối ưu duy nhất là $J^* = \{1, 2, 4\}$ với $S_1 \cup S_2 \cup S_4 = \{1, 2, 3, 4, 5, 6, 7\} = U$ và tổng chi phí tối thiểu $Z^* = c_1 + c_2 + c_4 = 7 + 5 + 4 = 16$. Trong khi đó, giải thuật tham lam Chvátal lần lượt chọn cột $4$ (tỷ số $4/4$), cột $5$ ($3/1$), cột $2$ ($5/1$) và cột $1$ ($7/1$) cho tổng chi phí $19$; sau bước tỉa cột dư thừa (loại bỏ cột $5$ do nút $4$ đã được cột $1$ phủ), nghiệm thu về $J = \{1, 2, 4\}$ đạt chi phí $16$.

\section{Tính chất NP-khó và hệ quả thuật toán}

\subsection{Chứng minh tính NP-đầy đủ}
Phiên bản quyết định của bài toán Phủ tập hợp (``Tồn tại hay không một tập chỉ số $J \subseteq \{1, \dots, n\}$ sao cho $\bigcup_{j \in J} S_j = U$ và $\sum_{j \in J} c_j \le K$?'') là một trong 21 bài toán \textbf{NP-đầy đủ} kinh điển của Richard Karp (1972):
\begin{enumerate}
    \item \textbf{Thuộc lớp NP:} Với một tập ứng viên $J$ và ngưỡng $K$ cho trước, ta kiểm tra điều kiện $\bigcup_{j \in J} S_j = U$ và tính tổng $\sum_{j \in J} c_j \le K$ trong thời gian đa thức $\mathcal{O}(m \cdot n)$.
    \item \textbf{Quy dẫn đa thức từ bài toán Phủ đỉnh (Vertex Cover):} Cho đồ thị vô hướng $G = (V, E)$ và số nguyên $k$ là một thực thể của bài toán Phủ đỉnh (đã biết là NP-đầy đủ). Ta biến đổi sang thực thể SCP trong thời gian đa thức bằng cách đặt tập vũ trụ $U = E$, và ứng với mỗi đỉnh $v \in V$ tạo một tập con $S_v = \{e \in E \mid e \text{ liên thuộc với } v\}$ với chi phí $c_v = 1$. Khi đó, $G$ có tập phủ đỉnh kích thước không quá $k$ khi và chỉ khi tồn tại tối đa $k$ tập con $S_v$ phủ toàn bộ $U$.
\end{enumerate}

\subsection{Khả năng xấp xỉ và định hướng thiết kế giải thuật}
Theo kết quả của Chvátal (1979), giải thuật tham lam đạt tỷ số xấp xỉ $H(d) \le 1 + \ln d$ với $d = \max_j |S_j|$, và theo Dinur cùng Steurer (2014), không tồn tại thuật toán thời gian đa thức nào đạt tỷ số xấp xỉ $(1 - \epsilon)\ln m$ với mọi $\epsilon > 0$ trừ khi $\text{P} = \text{NP}$. Vì vậy:
\begin{itemize}
    \item Trên các bộ dữ liệu kích thước nhỏ ($n \le 30$), nhóm cài đặt thuật toán chính xác Nhánh cận (\texttt{ExactSolver}) để tìm nghiệm tối ưu toàn cục làm mốc đối chiếu.
    \item Trên các bộ chuẩn lớn của OR-Library (như \texttt{scp4x}--\texttt{scp6x}, \texttt{scpa}--\texttt{scpd}, \texttt{scpnre}--\texttt{scpnrh}), nhóm sử dụng metaheuristic Simulated Annealing (\texttt{SASolver}) để tìm lời giải chất lượng cao trong ngân sách thời gian cho trước.
\end{itemize}

\section{Thuật toán chính xác: Nhánh cận (Branch and Bound)}

Lớp \texttt{ExactSolver} tìm nghiệm tối ưu toàn cục thông qua tìm kiếm theo chiều sâu (DFS) kết hợp giảm kích thước bài toán trước khi duyệt và cắt tỉa nhánh dựa trên cận dưới.

\subsection{Tiền xử lý giảm cỡ dữ liệu (Reduction)}
Trước khi bước vào cây tìm kiếm DFS, hàm \texttt{preprocess} thực hiện liên tiếp hai phép rút gọn bảo toàn tính tối ưu:
\begin{enumerate}
    \item \textbf{Loại bỏ cột rỗng và cột bị trội (Column Dominance):} Các cột không phủ hàng nào ($S_{j_1} = \emptyset$) bị loại bỏ. Ngoài ra, nếu tồn tại cặp cột $(j_1, j_2)$ thỏa mãn $c_{j_2} \le c_{j_1}$, $|S_{j_2}| \ge |S_{j_1}|$ và $S_{j_1} \subseteq S_{j_2}$, thì cột $j_1$ bị trội hoàn toàn bởi $j_2$ nên được đánh dấu loại bỏ khỏi không gian xét.
    \item \textbf{Cố định cột bắt buộc (Essential Columns):} Lặp lại quá trình quét các hàng chưa được phủ; nếu một hàng $i$ chỉ còn duy nhất một cột khả dụng $c$ bao phủ nó, thì cột $c$ bắt buộc phải thuộc mọi phép phủ hợp lệ. Thuật toán đưa $c$ vào danh sách \texttt{forced\_cols}, cộng $c_c$ vào chi phí nền \texttt{base\_cost}, cập nhật số lần phủ \texttt{cov\_cnt} của các hàng thuộc $S_c$ và đánh dấu không xét lại $c$ trong quá trình phân nhánh.
\end{enumerate}

\subsection{Chiến lược phân nhánh và cập nhật cận trên}
\begin{itemize}
    \item \textbf{Khởi tạo và cập nhật cận trên ($UB$):} Kỷ lục chi phí tốt nhất \texttt{best\_cost} được khởi tạo bằng $+\infty$. Mỗi khi nhánh duyệt phủ đủ toàn bộ $m$ hàng với tổng chi phí tích lũy $C_{\text{curr}} < \text{best\_cost}$, thuật toán cập nhật $\text{best\_cost} \leftarrow C_{\text{curr}}$ và ghi nhận danh sách cột tốt nhất \texttt{best\_cols}.
    \item \textbf{Chọn biến phân nhánh theo MRV:} Tại mỗi nút chưa phủ hết $m$ hàng, thuật toán tìm hàng chưa được phủ $i^*$ có số lượng cột khả dụng bao phủ nó ít nhất (\textit{Minimum Remaining Values} -- MRV) để thu hẹp hệ số phân nhánh. Các cột ứng viên chứa $i^*$ được sắp xếp theo chi phí $c_j$ tăng dần nhằm sớm tìm ra nghiệm hợp lệ có chi phí thấp, giúp siết chặt cận trên ngay từ những nhánh đầu tiên.
\end{itemize}

\subsection{Hàm ước lượng cận dưới ($LB$) và quy tắc cắt tỉa}
Tại nút hiện tại có chi phí tích lũy $C_{\text{curr}}$ và tập các hàng chưa được phủ là $U_{\text{rem}} = \{i \in U \mid \text{cov\_cnt}[i] = 0\}$, hàm \texttt{lower\_bound} tính chi phí tối thiểu ước lượng để phủ nốt $U_{\text{rem}}$ bằng cách phân bổ đều chi phí mỗi cột khả dụng cho toàn bộ kích thước $|S_j|$ của cột đó, sau đó làm tròn lên số nguyên gần nhất:
\[
LB = C_{\text{curr}} + \left\lceil \sum_{i \in U_{\text{rem}}} \min_{j : i \in S_j,\; j \text{ khả dụng}} \left( \frac{c_j}{|S_j|} \right) \right\rceil
\]
Nếu tồn tại một hàng $i \in U_{\text{rem}}$ không còn cột khả dụng nào phủ được, nhánh hiện tại vô nghiệm và bị loại ngay. Ngược lại, nếu $LB \ge \text{best\_cost}$ (hoặc vượt quá giới hạn thời gian \texttt{time\_limit\_sec}), thuật toán thực hiện cắt tỉa nhánh ngay lập tức.

\section{Thuật toán gần đúng: Simulated Annealing (SA)}

Lớp \texttt{SASolver} cài đặt thuật toán Luyện kim mô phỏng theo nguyên lý của Kirkpatrick và cộng sự (1983), duy trì tính hợp lệ của lời giải sau mỗi bước di chuyển thông qua cơ chế sửa chữa và loại bỏ dư thừa.

\subsection{Biểu diễn nghiệm và Khởi tạo tham lam Chvátal}
\begin{itemize}
    \item \textbf{Biểu diễn dữ liệu:} Lời giải hiện tại được lưu bằng mảng Boolean \texttt{in\_sol} kích thước $n + 1$ kết hợp mảng đếm \texttt{cov\_cnt} kích thước $m + 1$ ghi nhận số lượng cột đang được chọn bao phủ từng hàng $i \in U$. Cách tổ chức này cho phép kiểm tra tính dư thừa và cập nhật trạng thái trong thời gian $\mathcal{O}(|S_j|)$.
    \item \textbf{Khởi tạo tham lam và tỉa cột thừa:} Hàm \texttt{build\_initial} bắt đầu từ tập rỗng và ở mỗi bước chọn cột $j^*$ tối thiểu hóa tỷ số chi phí trên số hàng chưa phủ mới được bao phủ thêm:
    \[
    j^* = \arg\min_{j \notin J} \frac{c_j}{|S_j \cap U_{\text{rem}}|}
    \]
    Sau khi phủ toàn bộ $m$ hàng, hàm \texttt{remove\_redundant} duyệt qua các cột đang chọn theo thứ tự chi phí $c_j$ giảm dần; nếu mọi hàng $r \in S_j$ đều thỏa mãn $\text{cov\_cnt}[r] \ge 2$, cột $j$ được loại bỏ để giảm chi phí mà vẫn bảo toàn tính hợp lệ.
\end{itemize}

\subsection{Toán tử sinh láng giềng và phục hồi tính hợp lệ}
Tại mỗi bước lặp của chuỗi Markov, thuật toán tạo nghiệm ứng viên $J'$ từ nghiệm hiện tại $J$ qua ba giai đoạn:
\begin{enumerate}
    \item \textbf{Đảo bit ngẫu nhiên (Random Flip):} Chọn ngẫu nhiên $k \in \{1, 2\}$ cột (với xác suất $0{,}7$ đảo $1$ cột và xác suất $0{,}3$ đảo $2$ cột). Với mỗi cột $c$ được chọn ngẫu nhiên đồng đều trong $\{1, \dots, n\}$, thuật toán đảo trạng thái chọn của $c$ (nếu đang chọn thì loại bỏ, nếu chưa chọn thì thêm vào) đồng thời cập nhật gia tăng \texttt{next\_cost} và \texttt{next\_cov}.
    \item \textbf{Sửa chữa tham lam (Deterministic Greedy Repair):} Nếu việc loại bỏ cột khiến một số hàng bị mất bao phủ ($\text{next\_cov}[i] = 0$), hàm \texttt{repair} lần lượt duyệt các hàng chưa phủ theo chỉ số tăng dần và bổ sung cột $c$ phủ hàng $i$ có tỷ số chi phí trên số hàng mới phủ thêm $c_c / \max(1, \text{gain})$ nhỏ nhất.
    \item \textbf{Loại bỏ cột dư thừa (Redundancy Elimination):} Gọi lại thủ tục \texttt{remove\_redundant} theo thứ tự chi phí giảm dần nhằm loại bỏ các cột trở nên dư thừa sau bước đảo bit và sửa chữa.
\end{enumerate}

\subsection{Quy tắc chấp nhận Metropolis và Lịch làm nguội}
Gọi $\Delta = \text{Cost}(J') - \text{Cost}(J)$ là chênh lệch chi phí giữa nghiệm láng giềng khả thi $J'$ và nghiệm hiện tại $J$:
\begin{itemize}
    \item Nếu $\Delta \le 0$: Chấp nhận $J \leftarrow J'$ và cập nhật nghiệm tốt nhất toàn cục $J_{\text{best}}$ cùng nhật ký hội tụ (\texttt{conv\_csv\_file}) nếu $\text{Cost}(J') < \text{Cost}(J_{\text{best}})$.
    \item Nếu $\Delta > 0$: Chấp nhận nghiệm xấu hơn theo tiêu chuẩn Metropolis với xác suất $P(\Delta, T) = \exp(-\Delta / T)$ nhằm thoát khỏi cực trị địa phương.
\end{itemize}

\textbf{Tham số vận hành và tiêu chí dừng:}
\begin{itemize}
    \item \textbf{Nhiệt độ khởi tạo:} $T_0 = \max\big(1{,}0,\; 0{,}05 \times \text{Cost}(J_0)\big)$, tỷ lệ trực tiếp với chi phí nghiệm khởi tạo ban đầu.
    \item \textbf{Độ dài chuỗi Markov mỗi mức nhiệt:} $L = \max\big(20,\; \lfloor n / 5 \rfloor\big)$ bước lặp, tự động co giãn theo số lượng cột $n$ của thực thể.
    \item \textbf{Lịch giảm nhiệt độ:} Sau mỗi chuỗi $L$ bước lặp, nhiệt độ được cập nhật theo quy luật cấp số nhân có chặn dưới: $T \leftarrow \max(0{,}01,\; 0{,}995 \cdot T)$.
    \item \textbf{Tiêu chí dừng:} Thuật toán chạy liên tục và dừng khi thời gian thực thi chạm ngưỡng ngân sách thời gian \texttt{budget\_sec} được cấp phát.
\end{itemize}

\end{document}